\documentclass[letterpaper,12pt]{article}
\usepackage{graphicx}
\usepackage{pdflscape}
\usepackage{amsmath}
\usepackage{amsthm}
\usepackage{lscape}
\usepackage{enumerate}
\usepackage{float}
\usepackage{bm}
\usepackage{grffile}
\usepackage{relsize}
\usepackage{tablefootnote}
\usepackage{footnote}
\usepackage{sectsty}
\usepackage{xcolor}
\usepackage{caption}
\usepackage{color}
\usepackage{natbib}
\usepackage{multirow}
\usepackage{bbm}
\usepackage[T1]{fontenc}
\usepackage[multiple]{footmisc}
\usepackage{appendix}
\usepackage{times}
\usepackage{relsize}
\usepackage{newtxtext,newtxmath}
\usepackage{enumitem}

%%%%%%%%%%%%%%%%%%%%%%%%%%%%%%%%%%%%%%%%%%%%%%%%%%%%%%%%%%%%
%%%   Title Page Information
%%%%%%%%%%%%%%%%%%%%%%%%%%%%%%%%%%%%%%%%%%%%%%%%%%%%%%%%%%%%
\title{\vspace{-.5in} ECON712: Problem Set 3}
\author{Leandro Freylejer}
\date{Due date: April 15th before noon }

%%%%%%%%%%%%%%%%%%%%%%%%%%%%%%%%%%%%%%%%%%%%%%%%%%%%%%%%%%%%
%%%   Double and Single Space Commands
%%%%%%%%%%%%%%%%%%%%%%%%%%%%%%%%%%%%%%%%%%%%%%%%%%%%%%%%%%%%
\newlength{\defbaselineskip}
\setlength{\defbaselineskip}{\baselineskip}
\newcommand{\setlinespacing}[1]{\setlength{\baselineskip}{#1 \defbaselineskip}}
\newcommand{\doublespacing}{\setlength{\baselineskip}{1.3 \defbaselineskip}}
\newcommand{\singlespacing}{\setlength{\baselineskip}{1\defbaselineskip}}
\renewcommand{\thesubsection}{\arabic{subsection}}

%%%%%%%%%%%%%%%%%%%%%%%%%%%%%%%%%%%%%%%%%%%%%%%%%%%%%%%%%%%%
%%%   Set Margins Here
%%%%%%%%%%%%%%%%%%%%%%%%%%%%%%%%%%%%%%%%%%%%%%%%%%%%%%%%%%%%
\textwidth=6.5in
\textheight=9in
\oddsidemargin=0.10in
\evensidemargin=0.10in
\topmargin=-0.5in

%%%%%%%%%%%%%%%%%%%%%%%%%%%%%%%%%%%%%%%%%%%%%%%%%%%%%%%%%%%%%
%%%   Actual Document Begins Here
%%%%%%%%%%%%%%%%%%%%%%%%%%%%%%%%%%%%%%%%%%%%%%%%%%%%%%%%%%%%%
\begin{document}
\pagenumbering{arabic}

\maketitle
\vspace{-.3in}
\textbf{Instructions:} Answer all questions. To receive full marks, write your answers as completely and clearly as possible, providing all intermediate steps and stating any assumptions invoked.


%%%%%%%%%%%%%%%%%%%%%%%%%%%%%%%%%%%%%%%%%%%%%%%%%%%%%%%%%%%%%
\section*{Question I: Panel Data}
%%%%%%%%%%%%%%%%%%%%%%%%%%%%%%%%%%%%%%%%%%%%%%%%%%%%%%%%%%%%%

Consider a balanced panel dataset that observes $n$ units (e.g., firms) over $T$ time periods. Suppose the true data-generating process is given by the following \textbf{fixed effects model}:
\begin{equation}
    y_{it} = \bm{x}_{it}'\bm{\beta} + \alpha_i + \varepsilon_{it}, \qquad i = 1,\ldots,n, \quad t = 1,\ldots,T,
    \label{eq:true_model}
\end{equation}
where $\bm{x}_{it}$ is a $K \times 1$ vector of regressors, $\bm{\beta}$ is the $K \times 1$ coefficient vector of interest, $\alpha_i$ is a \textbf{time-invariant unit fixed effect} that may be correlated with $\bm{x}_{it}$, and $\varepsilon_{it}$ is an idiosyncratic error. Assume throughout that strict exogeneity holds: $\mathbb{E}[\varepsilon_{it} \mid \bm{x}_{i1},\ldots,\bm{x}_{iT},\alpha_i] = 0$.

\medskip
Stack all $NT$ observations into the $(NT \times 1)$ vector $\bm{y}$, the $(NT \times K)$ matrix $\bm{X}$, the $(NT \times 1)$ vector $\bm{\alpha}$ (where unit $i$'s fixed effect $\alpha_i$ is repeated $T$ times in the rows corresponding to that unit), and the $(NT \times 1)$ vector $\bm{\varepsilon}$. The true model in stacked form is:
\begin{equation}
    \bm{y} = \bm{X}\bm{\beta} + \bm{\alpha} + \bm{\varepsilon}.
    \label{eq:stacked}
\end{equation}
A researcher who does not account for $\alpha_i$ treats the composite term $\bm{u} \equiv \bm{\alpha} + \bm{\varepsilon}$ as the error and runs pooled OLS of $\bm{y}$ on $\bm{X}$.

\begin{enumerate}[label=(\alph*)]

\item (7 points) Derive the pooled OLS/MM estimator $\hat{\bm{\beta}}^{OLS}$ in matrix form. Using the Law of Large Numbers as $n\to\infty$ with $T$ fixed, derive $\mathrm{plim}\;\hat{\bm{\beta}}^{OLS}$ as a function of population parameters. 

\item (5 points) Under what conditions on the population parameters is $\hat{\bm{\beta}}^{OLS}$ inconsistent? Provide an economic interpretation and explain intuitively why pooled OLS fails in those cases.

\end{enumerate}


%%%%%%%%%%%%%%%%%%%%%%%%%%%%%%%%%%%%%%%%%%%%%%%%%%%%%%%%%%%%%
\section*{Question II: Instrumental Variables and Returns to Schooling}
%%%%%%%%%%%%%%%%%%%%%%%%%%%%%%%%%%%%%%%%%%%%%%%%%%%%%%%%%%%%%

This question replicates the instrumental variables analysis in Card (1995), following Cunningham's \textit{Causal Inference: The Mixtape} (Chapter 7, Table 57, p.\ 355). Card (1995) studies the returns to schooling using \textbf{proximity to a four-year college} (\texttt{nearc4}) as an instrument for completed years of education (\texttt{educ}), with log wages (\texttt{lwage}) as the outcome. The identifying idea is that individuals who grew up near a college face a lower cost of attending, inducing variation in educational attainment that is plausibly unrelated to unobserved determinants of wages.

\begin{enumerate}[label=(\alph*)]

%%% Part (a): FWL scatter plots + Wald estimate + weak instrument
\item (10 points) Use the Frisch--Waugh--Lovell (FWL) theorem to partial out the included controls and visualize the identifying variation in the data.

\begin{enumerate}[label=(\roman*)]
    \item Let $\bm{X}_c$ denote the matrix of included controls (all regressors except \texttt{educ} and \texttt{nearc4}: experience, its square, \texttt{black}, \texttt{south}, \texttt{smsa}, and region dummies, plus a constant). Regress each of \texttt{educ}, \texttt{nearc4}, and \texttt{lwage} separately on $\bm{X}_c$ and save the residuals $\tilde{d}$, $\tilde{z}$, and $\tilde{y}$, respectively. By FWL, these residuals represent the variation in each variable that is orthogonal to the included controls.

    \item Produce two scatter plots, each with a fitted regression line overlaid:
    \begin{itemize}
        \item \textbf{First stage:} $\tilde{d}$ (residualized education) on the vertical axis against $\tilde{z}$ (residualized \texttt{nearc4}) on the horizontal axis.
        \item \textbf{Reduced form:} $\tilde{y}$ (residualized log wage) on the vertical axis against $\tilde{z}$ on the horizontal axis.
    \end{itemize}
    Label your axes clearly and include the slope coefficient of each fitted line on the plot or in the caption.

    \item Using the slope coefficients from (ii), compute the Wald IV estimate as:
    \begin{equation}
        \hat{\beta}^{IV}_{Wald} = \frac{\text{slope of reduced form}}{\text{slope of first stage}} = \frac{\tilde{z}'\tilde{y}}{\tilde{z}'\tilde{d}}.
        \label{eq:wald}
    \end{equation}
    Report your estimate 

    \item Using the first-stage regression from (i), report the $F$-statistic on \texttt{nearc4} and comment on whether the instrument appears \textbf{weak}. 
\end{enumerate}

%%% Part (b): 2SLS replication of Table 57
\item (10 points) \textbf{2SLS estimation (replication of Table 57, p.\ 355).}
Replicate the OLS and 2SLS results reported in Table 57 of \textit{The Mixtape}. Interpret results and compare your 2SLS estimate to the IV-Wald estimate you obtained in part (a)


%%% Part (c): Instrument invalidity
\item (8 points) Describe the conditions under which \texttt{nearc4} would \emph{fail} to be a valid instrument in this setting. Provide at least two concrete, economically plausible channels through which college proximity could directly affect wages independently of years of schooling completed. For each channel, explain the direction of the resulting bias in the 2SLS estimate.

%%% Part (d): ATE vs LATE
\item (8 points) In words (no algebra required), explain what the 2SLS estimator identifies when the instrument is valid but \emph{heterogeneous treatment effects} are present across individuals.

\end{enumerate}

\end{document}
